% ECONOMICS_PROBLEM: {"id": "H43-447", "title": "Public-investment returns", "jel_code": "H43", "source_id": "OP-0252"}
% !TeX program = xelatex
\documentclass[11pt,a4paper]{article}
\usepackage[margin=23mm]{geometry}
\usepackage{fontspec}
\setmainfont{Latin Modern Roman}
\usepackage{amsmath,amssymb,mathtools}
\usepackage{xcolor,enumitem,booktabs,longtable,array,xurl,hyperref}
\definecolor{muted}{HTML}{53636C}
\hypersetup{colorlinks=true,urlcolor=blue,linkcolor=blue}
\setlength{\parindent}{0pt}
\setlength{\parskip}{5pt}
\newcommand{\R}{\mathbb R}
\newcommand{\E}{\mathbb E}
\newcommand{\N}{\mathbb N}
\newcommand{\ind}{\mathbf 1}
\newcommand{\Var}{\operatorname{Var}}
\newcommand{\Cov}{\operatorname{Cov}}
\newcommand{\supp}{\operatorname{supp}}
\DeclareMathOperator*{\argmin}{arg\,min}
\DeclareMathOperator*{\argmax}{arg\,max}
\newcommand{\entrymeta}[3]{{\sffamily\small\color{muted}\textbf{\texttt{#1}}\quad|\quad #2\par #3\par}}
\newcommand{\Problemref}[1]{\texttt{#1}}
\newcommand{\JELref}[2]{\texttt{#2} (JEL \texttt{#1})}
\begin{document}
\section*{Public-investment returns}
\textbf{Problem ID:} \texttt{H43-447}\quad \textbf{JEL:} \texttt{H43}\par
% BEGIN ECONOMICS STATEMENT
\label{problem:OP-0252}
\label{audited:ECON-097}
{\sffamily\small\color{muted}Original subject group: Fiscal policy and sovereign debt\par}
\label{definitions:problem:30007533}
\textbf{Audit classification:} Background; proposed extension. Official report summary is verified; project-specific net-value identification is not its stated open question. Verification concerns the cited version, not first priority or proof that this exact editorial specification remains unsolved on October 8, 2026.
\subsection*{Definitions and precise research target}
\paragraph{Editorial mathematical specification.} Fix a class of public infrastructure projects \(j\) with construction cost \(K_j\), maintenance stream \(M_{jt}\), completion delay \(d_j\), and independently measured service capacity. Let \(A_{jt}(1)\) and \(A_{jt}(0)\) denote private-sector productivity with and without project implementation. Let \(S_{jt}(A)\) be the gross social value of output and access at productivity level \(A\), net of congestion and environmental damages; its valuation assumptions must be declared. At discount rate \(r>0\), define
\[
V_j=E\sum_{t\ge d_j}(1+r)^{-t}
[S_{jt}(A_{jt}(1))-S_{jt}(A_{jt}(0))-M_{jt}]
-K_j-C_j^F,
\]
where \(C_j^F\) is the welfare cost of the specified financing scheme, including displaced expenditures or distortionary taxes.
Identify conditional expected \(V_j\) and mean productivity gains; a distribution of realized individual gains requires additional joint-outcome restrictions. Use credible project assignment, timing, or funding instruments to identify the stated mean effects. Address endogenous project location, anticipation, spillovers to untreated areas, and selective completion. Separate physical infrastructure productivity from utilization and administrative efficiency, and quantify the consequences of delays and cost overruns. Establish which measured project characteristics predict positive net value out of sample. A country-level spending regression cannot by itself identify project returns or financing costs. This is a scoped project-evaluation problem; the general possibility that productive public investment benefits growth is already supported by substantial research.

Express \(S,M,K,C^F\) in a common date-zero consumption numeraire; a utility loss cannot be subtracted directly from a monetary project return. If \(C^F\) includes repayment of the capital outlay already counted in \(K\), remove that repayment component and retain only additional distortion or opportunity cost. Maintenance is incremental relative to the no-project regime, and construction costs may be a dated stream rather than all paid at date zero. Completion delay may be random; replace the summation by indicators \(1\{t\ge d_j\}\) and require absolute integrability of discounted benefits and costs. Productivity \(A\) is a specified production residual and does not measure access benefits, which must be valued separately. Even a randomized project reveals its marginal potential-outcome distributions, not the distribution of realized individual gains, without a joint-outcome restriction.
\subsection*{Variants and assumption changes}
The following are \emph{editorial variants}, not additional published open conjectures.
\begin{enumerate}
\item \emph{Mean net present value.} Under randomized assignment or a valid local assignment design identify \(E[V_j]\) for the affected population, with explicit spillover exposure and incremental maintenance. A positive mean does not imply every project has positive value.
\item \emph{Distributional and transport bounds.} Use baseline characteristics to bound the fraction of projects with positive net value and extend estimates to untested regions under quantified transport restrictions. Avoid asserting joint potential-outcome identification from marginal treatment means.
\end{enumerate}
\subsection*{References and source audit}
\begin{itemize}
\item International Monetary Fund, Fiscal Affairs Department. \emph{Fiscal Monitor, October 2025: Spending Smarter—How Efficient and Well-Allocated Public Spending Can Boost Economic Growth}. 2025. Locator: Official Fiscal Monitor October 2025 publisher summary; Chapter1 description; indexed chapter DOI10.5089/9798229024938.089.CH001. Official publisher summary checked; direct chapter access was forbidden and full report fetch failed. \url{https://www.imf.org/en/publications/fm/issues/2025/10/07/fiscal-monitor-october-2025}.
\end{itemize}
Official report summary is verified; project-specific net-value identification is not its stated open question.
\begin{enumerate}\raggedright
\item\hypertarget{definitions:source:30007533:1}{} International Monetary Fund, Fiscal Affairs Department. 2025. \emph{Fiscal Monitor, October 2025: Spending Smarter—How Efficient and Well-Allocated Public Spending Can Boost Economic Growth}. Official Fiscal Monitor October2025 publisher summary; Chapter1 description; indexed chapter DOI10.5089/9798229024938.089.CH001.
\url{https://www.imf.org/en/publications/fm/issues/2025/10/07/fiscal-monitor-october-2025}.
DOI: \url{https://doi.org/10.5089/9798229024938.089}.
\end{enumerate}

\subsection*{Integrated refinement: ECON-097}
{\sffamily\small\color{muted}Network-wide social returns to transport investment\par}
This refinement adopts the additional source conventions
(page~\pageref{audited:conventions}), subject to its local restrictions.
\textbf{Network-wide transport returns and accounting boundaries.}
For a capacity/travel-cost project $z$, allow routes, residence, production and
trade to adjust in a selected equilibrium. For a fixed welfare population use
\[
NPV(z)=\sum_{t=0}^{T}\delta^t
[\Delta B_t^{\rm eq}(z)-\Delta D_t^{\rm omitted}(z)-C_t^{\rm omitted}(z)].
\]
The monetary equilibrium benefit includes labor and ownership income once;
subtract only damages and real project costs absent from it. An equivalent
gross-benefit convention instead includes all real costs exactly once.
Specify life-cycle construction, maintenance, delay, congestion, safety and
environmental costs with a common numeraire. Identify network-wide and
distributional gains under explicit spillover and location-choice assumptions
rather than summing direct user effects or adding profits already in household
income.
\subsection*{Supplied source audit for ECON-097}
Local [S1], [S2], etc. in this audit and the references immediately below
belong to ECON-097, independently of the original entry's reference numbering.
[S1] and [S2] support cost-benefit data and network/spatial transport returns. This NPV specification is editorial, with profit and capitalization double-counting removed.
\subsection*{References for integrated refinement ECON-097}
{\footnotesize\raggedright
\par\noindent\textbf{[S1]} Marco Gonzalez-Navarro, Román David Zárate, Rémi Jedwab, and Nick Tsivanidis (2023). \emph{Land Transport Infrastructure}. VoxDevLit 9(1), December 2023.\par \textit{Verified locator / source role:} Primary publisher page checked for review identity, authors, version and background. The specific published agenda is corroborated in [S2]; the exact formal target is editorial.\par \url{https://voxdev.org/voxdevlit/land-transport-infrastructure}\par\smallskip
\par\noindent\textbf{[S2]} Oliver Hanney / VoxDev (living compilation). \emph{Evidence gaps: Key unanswered questions in development economics}. Accessed 8 October 2026. The page currently covers 20 topics; the ``16-topics URL is a historical slug.\par \textit{Verified locator / source role:} Topic heading: Land Transport Infrastructure. Supports the broad research agenda; this is not a verbatim quotation or an attribution of the displayed mathematical target.\par \url{https://voxdev.org/topic/evidence-gaps-key-unanswered-questions-across-16-topics-development-economics}\par\smallskip
}

\subsection*{Common conventions: Source mathematical conventions}

These conventions apply to each entry unless explicitly replaced.

\textbf{Models and identification.} A primitive model \(m\) specifies all probability laws, feasible choices, information, equilibrium and measurement rules used by its target. The admissible class \(\mathcal M\) is fixed independently of outcomes. If \(P_O(m)\) is its observable probability law and \(\tau(m)\) the target, its identified set at observable law \(P\) is
\[
 \mathcal I_\tau(P)=\{\tau(m):m\in\mathcal M,\ P_O(m)=P\}.
\]
Point identification means that this set is a singleton. An empty set signals incompatibility of data and assumptions. Coordinate bounds need not describe the joint set. Bounds are sharp only if every retained target is attainable or arbitrarily approachable by compatible models. Finite bounds require explicit restrictions on unobserved quantities; unrestricted confounding, hidden wealth, extrapolation or arbitrary equilibrium selection can yield uninformative sets.

\textbf{Causal interventions.} A potential outcome \(Y_i(p)\) is the outcome under a complete policy regime \(p\), including timing, financing, information and market spillovers. The target population and units are fixed before treatment. With interference, use the complete assignment vector or a justified exposure mapping; individual no-interference notation alone cannot identify an economy-wide intervention. A difference of mean potential outcomes does not require the cross-world joint distribution. Conditioning on post-treatment migration, survival, enrollment or employment changes the estimand and needs separate justification.

\textbf{Statistical statements.} An estimator, confidence set, forecast or test specifies a sampling law, model class and loss/coverage criterion. Uniform \((1-\alpha)\)-coverage means \(\inf_{m\in\mathcal M}\Pr_m\{\tau(m)\in C_n\}\geq1-\alpha\), or an explicitly asymptotic version. The whole parameter space trivially has coverage, so informativeness requires a stated width, risk or power objective. A forecasting improvement is relative to a named benchmark, information set and evaluation population.

\textbf{Equilibrium and optimization.} An equilibrium comprises feasible plans, optimality, beliefs where needed, budget constraints and all market/resource clearing. The primitives or a stated benchmark define each correspondence \(\mathcal E(p;m)\). Nonemptiness is assumed or proved, never inferred just from compact policy sets. A selected-equilibrium target uses a declared selection rule; a robust target quantifies over all permitted equilibria. Use suprema or infima unless attainment is established. An \(\varepsilon\)-optimal policy is feasible and within \(\varepsilon\) of the finite optimum value. Report an empty feasible set separately.

\textbf{Welfare and accounting.} Normative utility, interpersonal normalization, social weights, numeraire, discounting and the use of public receipts are primitives. Behavioral utility need not coincide with the welfare criterion. Household welfare includes ownership income and rents once; adding profits again to compensating variation generally double-counts them. A monetary equivalent is defined by an explicit transfer and price convention, with monotonicity and range conditions. Average effects, mechanism contributions, transfers and welfare losses are different targets. Infinite sums require convergence and dynamic feasible plans require terminal or no-Ponzi conditions.

\textbf{Source versions.} A working-paper date, revision date and journal publication year are distinguished. A DOI for a journal version is not automatically the DOI of an earlier manuscript. Locators refer to the stated version and printed pages where specified. A metadata-only check does not verify a claim about a full-text conclusion. Every entry's source subsection narrows the provenance claim accordingly.

\subsection*{Common conventions: Additional-source status and mathematical conventions}

\label{audited:conventions}

Of the 100 supplied ECON entries, 65 distinct problems appear here; 32 close
matches refine earlier entries, and three duplicate questions map to existing
entries. Appendix~\ref{app:audited-crosswalk} lists every destination.
The attachment's P/R/S/U distinctions and source audits are retained. Its
precise mathematical formalizations are editorial unless the individual audit
says otherwise. Candidate subsidy targets ECON-004--006 retain their U flags;
the resolved approval-core question ECON-007 updates OP-0202 above.
The following supplied conventions govern both this part and the attributed
ECON refinements elsewhere in the catalogue, subject to local restrictions.

\textbf{Parameterized questions.} A problem family lists inputs that must be fixed before an instance is evaluated: population, horizon, policies, information, data law, model class and welfare weights. These are required choices, not quantities the citations are assumed to specify. Undefined applications should not be treated as identified estimands.

\textbf{Indices and integrability.} Sets of agents, firms and regions are finite unless a population measure is explicitly used. \(i,j\) index units, \(r\) regions, \(t\) dates and \(h\) horizons. \(\mathbb E\) and \(\Pr\) refer to the declared population and law; \(\mathbf1\{A\}\) is an indicator. Every displayed expectation and discounted sum needs existence in the stated domain. Infinite horizons use a discount in \((0,1)\) and a convergence condition; finite horizons may use discount one.

\textbf{Optimization and existence.} A displayed \(\max\), \(\min\), or \(\arg\max\) is conditional on attainment. For finite feasible menus this is immediate; general spaces need continuity and compactness/coercivity or another existence argument. Without attainment, the target is the supremum/infimum and arbitrarily near-optimal feasible choices. \(\inf\varnothing=+\infty\). Empty feasibility and an unidentified target are different issues.

\textbf{Potential outcomes.} \(Y_i(z)\) is the outcome under a specified intervention, not the observed conditional mean among people choosing \(z\). In binary comparisons \(\Delta Y_i=Y_i(1)-Y_i(0)\). Specify assignment, treatment versions, compliance, information and observation horizon. Interference is allowed under a full policy vector; compressed exposure notation needs a justified mapping. No interference, consistency and overlap are substantive assumptions, not automatic consequences of notation.

\textbf{Populations and observation.} Fix baseline eligibility and population weights. Conditioning on post-policy employment, survival, relocation or program uptake can change the population being compared. Death is not ordinary loss to follow-up; outcomes truncated by death need a composite convention, joint survival target, or explicitly defined survivor stratum. Fertility-changing interventions need a population functional rather than an unexplained fixed descendant cohort. Measurement and missingness models must be supplied.

\textbf{Identification.} A structural model \(m\in\mathcal M\) implies observable law \(P_m\) and target \(\theta(m)\). Identification means
\[
 P_{m_1}=P_{m_2}\ \Longrightarrow\ \theta(m_1)=\theta(m_2)
 \quad\text{for all }m_1,m_2\in\mathcal M .
\]
Without identification, the target is
\[
 \Theta(P;\mathcal M)=\{\theta(m):m\in\mathcal M,\ P_m=P\}.
\]
Writing \(\theta(P)\) before establishing identification can hide incompatible latent models. Confidence sets additionally need a sampling law, confidence level, asymptotic sequence and uniformity class. Point identification, coverage and informative precision are separate properties.

\textbf{Mediation.} Total effects of randomized assignment are distinct from cross-world natural indirect effects. Belief/effort-channel effects require a meaningful mediator intervention and additional measurement, exclusion or mediation assumptions. Randomization of the initial policy alone does not supply those assumptions. Report bounds or interventional alternatives when the stronger target is unsupported.

\textbf{Equilibrium.} A counterfactual \(W(z)\), \(p(z)\), or \(\mu_t^z\) requires individual optimization, feasibility, government/resource budgets, market clearing and state-transition laws. State equilibrium existence and selection, or report ranges over compatible equilibria. Initial-state integration includes subsequent endogenous transitions; current-state integration must use the policy-dependent distribution. Reduced-form invariance after policy is not assumed.

\textbf{Welfare accounts.} Scores, utility, health, dollars and ecological quantities require explicit conversion before addition. Fees, taxes, subsidies, wages and extortion payments are transfers between parties; distributional weights can make them welfare relevant, but their face value is not automatically a resource loss. State ownership, geographic boundaries, foreign-income weights and fiscal finance. Do not add profits to household welfare already including dividends, or count land capitalization and the underlying benefit twice. A benefit-cost expression must distinguish gross benefits from net welfare and subtract every real cost exactly once.

\textbf{Derivatives and risk.} Log derivatives require positive arguments; local responses require admissible perturbations and specified equilibrium branches. Differentiation under expectations needs regularity/domination, and boundaries may allow only one-sided derivatives. Risk uses a specified law; ambiguity uses a declared nonempty law set on a common history space. Adaptive policies must be nonanticipative. A robust criterion is an editorial choice unless attributed explicitly.

\textbf{Scope overlaps.} Allocation, collective choice, contracts, household bargaining and coordinated policy overlap economics with optimization and game theory. They are retained because they appeared in the original catalogue; no claim of a strictly game-theory-free selection is made.
% END ECONOMICS STATEMENT
\end{document}
